\documentclass[10pt,letterpaper,oneside]{article}
\usepackage[T1]{fontenc}
\usepackage[utf8]{inputenc}
\usepackage{lmodern}
\usepackage{geometry}
\geometry{margin=0.84in,headheight=21pt,footskip=28pt}
\usepackage{microtype}
\usepackage{xcolor}
\definecolor{ink}{HTML}{15253D}
\definecolor{accent}{HTML}{086A84}
\definecolor{pale}{HTML}{EEF5F8}
\definecolor{muted}{HTML}{4E6474}
\definecolor{grayline}{HTML}{D7E2E8}
\usepackage{graphicx}
\usepackage{fancyhdr}
\usepackage{titlesec}
\usepackage{enumitem}
\usepackage{array,tabularx,longtable,booktabs}
\usepackage[most]{tcolorbox}
\usepackage{xurl}
\usepackage{hyperref}
\usepackage{listings}
\usepackage{amsmath,amssymb}
\usepackage{tikz}
\usetikzlibrary{arrows.meta,positioning,shapes.geometric,calc}
\hypersetup{colorlinks=true,linkcolor=accent,urlcolor=accent,pdftitle={Hands-On Security Engineering Laboratory Handbook},pdfauthor={Independent technical engineering study companion}}
\renewcommand{\familydefault}{\sfdefault}
\linespread{1.07}
\setlength{\parindent}{0pt}
\setlength{\parskip}{5.5pt}
\emergencystretch=2em
\setlist[itemize]{leftmargin=1.5em,itemsep=2.5pt,topsep=3pt}
\setlist[enumerate]{leftmargin=1.7em,itemsep=3pt,topsep=3pt}
\titleformat{\section}{\Large\bfseries\color{ink}}{\thesection}{0.65em}{}[\vspace{-2pt}{\color{accent}\titlerule}]
\titleformat{\subsection}{\large\bfseries\color{accent}}{\thesubsection}{0.6em}{}
\titleformat{\subsubsection}{\bfseries\color{ink}}{\thesubsubsection}{0.6em}{}
\titlespacing*{\section}{0pt}{16pt}{9pt}
\titlespacing*{\subsection}{0pt}{11pt}{6pt}
\pagestyle{fancy}\fancyhf{}\fancyhead[L]{\footnotesize\color{muted} COMPUTER \& INFORMATION SECURITY ENGINEERING}\fancyhead[R]{\footnotesize\color{accent} SUITE / 7}\fancyfoot[L]{\footnotesize\color{muted} Independent study companion \quad|\quad 2026}\fancyfoot[R]{\footnotesize\color{muted}\thepage}\renewcommand{\headrulewidth}{0.4pt}\renewcommand{\headrule}{\hbox to\headwidth{\color{grayline}\leaders\hrule height \headrulewidth\hfill}}
\newtcolorbox{sourcebox}{enhanced,breakable,boxrule=0pt,leftrule=3pt,colframe=accent,colback=pale,arc=0pt,left=9pt,right=7pt,top=7pt,bottom=7pt}
\newtcolorbox{warningbox}{enhanced,breakable,boxrule=0.5pt,colframe=grayline,colback=white,arc=1pt,left=9pt,right=9pt,top=7pt,bottom=7pt}
\newcommand{\handbookref}[1]{\textsf{\color{accent}\textbf{[H:#1]}}}
\newcommand{\modernext}{\textsf{\color{muted}\textbf{[E]}}}
\newcommand{\designed}{\textsf{\color{accent}\textbf{[D]}}}
\newcommand{\deliverable}[1]{\textbf{Deliverable:} #1}
\lstdefinestyle{lab}{basicstyle=\ttfamily\scriptsize,keywordstyle=\color{accent}\bfseries,commentstyle=\color{muted},numbers=left,numberstyle=\tiny\color{muted},numbersep=7pt,breaklines=true,showstringspaces=false,frame=single,rulecolor=\color{grayline},backgroundcolor=\color{pale},xleftmargin=1em,framexleftmargin=0.6em,tabsize=2,columns=fullflexible}
\lstset{style=lab}
\renewcommand{\arraystretch}{1.21}
\newcommand{\makereportcover}[4]{%
\thispagestyle{empty}\vspace*{1.15in}
{\small\bfseries\color{accent} COMPUTER AND INFORMATION SECURITY ENGINEERING}\par\vspace{13pt}
{\Huge\bfseries\color{ink} #1\par}\vspace{15pt}
{\Large\color{muted} #2\par}\vspace{18pt}
\color{accent}\rule{0.94\textwidth}{2.1pt}\par\vspace{20pt}
{\normalsize\bfseries Document #3 of 8 \quad\textbullet\quad 16 October 2026\par}
\vspace{0.35in}
\begin{sourcebox}\textbf{Orientation.} #4\end{sourcebox}
\vfill
{\small\color{muted} Independent, original engineering study and practice materials grounded in\par
\emph{Computer and Information Security Handbook}, 4th ed., edited by John R. Vacca,\par
Elsevier/Morgan Kaufmann, 2025, Volumes I--II. Not a reproduction or official standard.}
\newpage\setcounter{page}{1}\tableofcontents\newpage
}
\begin{document}
\makereportcover{Hands-On Security Engineering Laboratory Handbook}{Controlled labs, reproducible commands, evidence bundles and acceptance tests}{7}{Eight safe, nonproduction learning labs with explicit objectives, prerequisites, commands, rollback, expected evidence and grading criteria.}

\section{Laboratory charter and reproducibility}
The handbook routinely pairs chapters with review questions and hands-on projects; the labs here are newly authored to bring the concepts together. \handbookref{5,10--12,31,37,41--46,75--82} \designed{}
\begin{warningbox}\textbf{Authorization and safety:} perform only on local VMs, containers, loopback interfaces, test repositories, or systems where written scope authorizes the activity. Use synthetic nonpersonal test data. Do not scan third parties or deploy weak example settings into production. Always record the environment, expected outputs and rollback steps.\end{warningbox}
\subsection{Minimum lab environment}
Linux VM with non-root account, Python 3.10+, \texttt{curl}, \texttt{openssl}, standard networking tools; optionally \texttt{tcpdump}, Docker/Podman and a private test repository. Install optional tools from trusted OS/vendor sources after reviewing their version and licenses. No privileged network scanning is required to complete core labs.
\subsection{Evidence manifest schema}\designed{}
\begin{lstlisting}[caption={Example evidence bundle manifest (illustrative JSON)}]
{
  "lab": "LAB-02", "environment": "local-vm",
  "time_utc": "2026-10-16T17:00:00Z",
  "operator": "student-alias", "scope": "127.0.0.1 only",
  "tool_versions": {"python": "record actual version"},
  "artifacts": [{"name": "requests.txt", "sha256": "record digest"}],
  "result": "pass", "limitations": []
}
\end{lstlisting}
Do not place credentials, production PII or private keys in evidence bundles. Record timestamps and software versions at runtime rather than copying the example values.

\section{LAB-01: Host and privilege inventory}
\textbf{Outcome:} build a baseline inventory from a machine you administer; classify configuration drift. \handbookref{11,12,78,80}
\begin{enumerate}
\item Capture OS/kernel information, network listening ports and enabled services.
\item Enumerate current account and privilege context without accessing another user's secrets.
\item Identify unused services and patch ownership; recommend rather than automatically changing settings.
\item Create baseline controls \texttt{HOST-01} (supported OS), \texttt{HOST-02} (limited administrative access), \texttt{HOST-03} (approved listeners).
\end{enumerate}
\begin{lstlisting}[language=bash,caption={Read-only Linux host reconnaissance}]
cat /etc/os-release
uname -r
id
ss -lnt
command -v systemctl && systemctl --type=service --state=running
\end{lstlisting}
\textbf{Evidence:} \texttt{host-baseline.txt}, exceptions list and three verification tests. \textbf{Rollback:} none (read-only). \textbf{Pass:} each discovered service has owner and business purpose. \designed{}

\section{LAB-02: Loopback HTTP observation and authorization}
\textbf{Outcome:} observe a synthetic HTTP request and demonstrate object-level authorization. \handbookref{10,75,79} \modernext{}
\begin{lstlisting}[language=bash,caption={Local-only HTTP exercise}]
mkdir -p /tmp/cise-http-lab
printf '%s\n' 'sample non-sensitive data' > /tmp/cise-http-lab/index.html
cd /tmp/cise-http-lab
python3 -m http.server 8765 --bind 127.0.0.1
# From a second terminal:
curl -i http://127.0.0.1:8765/
\end{lstlisting}
\textbf{Exercise:} record status line, headers, method and body. Explain why this simple test server provides \emph{no authentication or authorization} and is not a production service. Replace it in the design exercise with a route requiring a tenant/object permission check; define successful and cross-tenant negative tests. \textbf{Rollback:} stop the server with Ctrl-C and delete the temporary directory. \designed{}

\section{LAB-03: Packet inspection within owned scope}
\textbf{Outcome:} distinguish packet evidence from a high-level log. \handbookref{5,44,75}
\begin{lstlisting}[language=bash,caption={Optional local-only capture}]
# Privileges for packet capture vary by platform.
sudo tcpdump -i lo -nn -s 0 -c 20 'tcp port 8765'
# Generate traffic with the curl request from LAB-02.
\end{lstlisting}
If the environment lacks permission for capture, use its approved loopback capture tool or a synthetic training trace. \textbf{Expected evidence:} observed source/destination addresses, port, sequence and timestamps; note when payloads are unavailable, encrypted or truncated. \textbf{Privacy:} never capture real credentials. \textbf{Rollback:} terminate capture; protect or delete capture artifacts under retention policy. \designed{}

\section{LAB-04: TLS certificate inspection}
\textbf{Outcome:} analyze identity binding and certificate validity without creating a false trust exception. \handbookref{3,49} \modernext{}
\begin{lstlisting}[language=bash,caption={Local trust-store-backed TLS test}]
# Use a trusted public test service only for normal client inspection.
# Substitute an approved organization-owned HTTPS endpoint if desired.
openssl version
openssl s_client -connect example.com:443 -servername example.com \
    -verify_return_error </dev/null
\end{lstlisting}
\textbf{Tasks:} identify trust anchor, server name, validity interval and any verification error; compare a correct SNI name with a mismatched name \emph{in a controlled test}. \textbf{Evidence:} sanitized chain summary with private data excluded. \textbf{Pass:} explain that a certificate chain alone does not authorize a user's document access.

\section{LAB-05: Local secure-code negative tests}
\textbf{Outcome:} encode object authorization as an automatically tested invariant. \handbookref{10,31,79} \modernext{}
\begin{lstlisting}[language=Python,caption={Save as test\_access.py; stdlib unittest}]
import unittest

def can_read(user_tenant, object_tenant, permissions):
    return user_tenant == object_tenant and "read" in permissions

class AccessTests(unittest.TestCase):
    def test_owner(self):
        self.assertTrue(can_read("a", "a", {"read"}))
    def test_cross_tenant(self):
        self.assertFalse(can_read("a", "b", {"read"}))
    def test_no_permission(self):
        self.assertFalse(can_read("a", "a", set()))

if __name__ == "__main__":
    unittest.main()
\end{lstlisting}
\begin{lstlisting}[language=bash,caption={Run tests and preserve exit status}]
python3 -m unittest -v test_access.py
\end{lstlisting}
\textbf{Extension:} add delegated access, inactive user and confused-deputy cases. \textbf{Pass:} all tests pass; explain which production behaviors remain untested (persistence, caching and concurrency). \designed{}

\section{LAB-06: Controlled vulnerability triage and code scanning}
\textbf{Outcome:} create a defensible triage record without conflating severity with exploitability. \handbookref{31,32,33} \modernext{}
\begin{enumerate}
\item Create a private, synthetic test repository with a known insecure function and a repaired version.
\item Run a permitted local linter or repository security analysis; record exact commit, tool and rule.
\item Classify issue and affected asset; explain preconditions and possible compensating controls.
\item Implement a fix; rerun analysis and unit tests on the candidate change and merged deployment where applicable.
\item Complete fields: finding ID, CWE, severity basis, confidence, technical impact, disposition, evidence, owner and verification.
\end{enumerate}
\subsection{Example security CI workflow}\modernext{}
The following checks out source and runs local unit tests; it is a minimal teaching gate, \emph{not} a complete code scanning deployment. Additional SAST/SCA/secret scanning jobs require maintained tools and organization-specific licensing/configuration.
\begin{lstlisting}[caption={.github/workflows/security-smoke.yml}]
name: Security smoke
on:
  pull_request:
  push:
    branches: [main]
jobs:
  test:
    runs-on: ubuntu-latest
    permissions:
      contents: read
    steps:
      - uses: actions/checkout@v4
      - uses: actions/setup-python@v5
        with:
          python-version: '3.12'
      - run: python -m unittest discover -v
\end{lstlisting}
\textbf{Production note:} pin action versions according to current organizational standards, preferably verified immutable commit SHAs; these illustrative major version tags are not a guarantee of newest releases. Require check results on the intended ref. \designed{}

\section{LAB-07: Backup and recovery proof}
\textbf{Outcome:} measure actual restore time and compare against an approved objective. \handbookref{9,37,38} \modernext{}
\begin{lstlisting}[language=bash,caption={Reversible synthetic-data restore}]
mkdir -p /tmp/cise-restore/{original,backup,restored}
printf '%s\n' 'known-good-record' > /tmp/cise-restore/original/item.txt
cp -a /tmp/cise-restore/original/. /tmp/cise-restore/backup/
sha256sum /tmp/cise-restore/original/item.txt
cp -a /tmp/cise-restore/backup/. /tmp/cise-restore/restored/
sha256sum /tmp/cise-restore/restored/item.txt
cmp /tmp/cise-restore/original/item.txt /tmp/cise-restore/restored/item.txt
\end{lstlisting}
\textbf{Evidence:} matching hashes, measured RTO, simulated RPO and recovery narrative. \textbf{Caveat:} this demonstrates file-copy correctness, not a real disaster recovery capability. \textbf{Rollback:} delete \texttt{/tmp/cise-restore} after retaining approved synthetic evidence.

\section{LAB-08: Forensic evidence and incident tabletop}
\textbf{Outcome:} preserve a small synthetic log and defend a time-ordered investigative conclusion. \handbookref{41--46} \modernext{}
\begin{lstlisting}[language=bash,caption={Hash a non-sensitive training evidence file}]
printf '%s\n' \
 '2026-10-16T09:02:00Z,login,privileged' \
 '2026-10-16T09:04:00Z,storage-list,privileged' \
 > /tmp/cise-training-evidence.csv
sha256sum /tmp/cise-training-evidence.csv
\end{lstlisting}
\textbf{Tabletop injects:} token detected outside normal region; object storage list operation; unexplained large outbound traffic; production owner requests immediate shutdown. Define a decision log, competing hypotheses, containment authorization and acquisition plan. \textbf{Pass:} separate confirmed login from suspected exfiltration, preserve custody, and specify restoration verification. \textbf{Rollback:} remove synthetic file after recording approved evidence.

\section{Lab completion matrix}
\par\medskip\noindent\begin{longtable}{@{}p{0.11\textwidth}p{0.33\textwidth}p{0.24\textwidth}p{0.22\textwidth}@{}}\toprule
Lab & Evidence artifact & Core competency & Pass criterion \\ \midrule\endhead
01 & Host baseline + exception register & Platform controls & Inventory/ownership complete\\
02 & Sanitized HTTP transcript & Protocol/authorization & Distinguishes authN from authZ\\
03 & Scoped packet notes & Network observation & States capture blind spots\\
04 & Certificate validation report & PKI & Explains trust chain and hostname\\
05 & Unit test log & Secure coding & Negative tests pass\\
06 & Triage finding + PR test run & Vulnerability mgmt & Provenance and fix traceable\\
07 & Restore verification record & Resilience & Actual restore measured\\
08 & Case timeline and custody sheet & DFIR & Evidence vs hypothesis separated\\
\bottomrule\end{longtable}\par\medskip

\section{Laboratory grading and expansion}
Each lab is assessed on safe scope (20\%), reproducibility (20\%), observed evidence (25\%), technical explanation (20\%) and cleanup (15\%). Pass: 80\% or higher \emph{and} no unsafe deviation from authorization. Suggested extension: join labs 02, 05 and 06 into one secure-API project and run a change-control walkthrough. \designed{}

\section*{Handbook chapter alignment}\addcontentsline{toc}{section}{Handbook chapter alignment}
These chapters ground the domain scope. The document offers original synthesis and technical procedures, not a chapter-by-chapter reproduction.
\par\medskip\noindent\begin{longtable}{@{}p{0.085\textwidth}p{0.68\textwidth}p{0.16\textwidth}@{}}\toprule Chapter & Handbook topic & Combined PDF p. \\ \midrule\endhead3 & A Cryptography Primer & 59 \\
5 & Detecting System Intrusions & 91 \\
9 & Fault Tolerance and Resilience in Cloud Computing Environments & 191 \\
10 & Securing Web Applications, Services, and Servers & 211 \\
11 & UNIX and Linux Security & 233 \\
12 & Eliminating the Security Weakness of Linux and UNIX Operating Systems & 253 \\
31 & What Is Vulnerability Assessment? & 573 \\
32 & Introduction to Vulnerabilities Below the Operating System & 585 \\
33 & Security Metrics: An Introduction and Literature Review & 595 \\
37 & Disaster Recovery & 651 \\
38 & Disaster Recovery Plans for Small and Medium Business (SMB) & 661 \\
41 & Cyber Forensics & 687 \\
42 & Cyber Forensics and Incidence Response & 717 \\
43 & Securing e-Discovery & 743 \\
44 & Network Forensics & 769 \\
45 & Microsoft Office and Metadata Forensics: A Deeper Dive & 785 \\
46 & Hard Drive Imaging & 795 \\
49 & Public Key Infrastructure & 841 \\
74 & Intrusion Prevention and Detection Systems & 1289 \\
75 & Transmission Control Protocol/Internet Protocol Packet Analysis & 1303 \\
76 & Firewalls & 1317 \\
77 & Penetration Testing & 1337 \\
78 & System Security & 1345 \\
79 & Access Controls & 1353 \\
80 & Endpoint Security & 1361 \\
81 & Assessments and Audits & 1369 \\
82 & Fundamentals of Cryptography & 1377 \\
\bottomrule\end{longtable}\par\medskip

\section*{Source and reference register}
\addcontentsline{toc}{section}{Source and reference register}
\begin{description}[leftmargin=1.6em,style=nextline]
\item[H] John R. Vacca (ed.), \emph{Computer and Information Security Handbook}, 4th ed., Volumes I--II, Elsevier/Morgan Kaufmann, 2025. Chapter identifiers in this document point to this book.
\item[N1] NIST, \emph{Cybersecurity Framework 2.0} (2024), \url{https://www.nist.gov/cyberframework}.
\item[N2] NIST SP 800-218, \emph{Secure Software Development Framework v1.1} (2022), \url{https://csrc.nist.gov/pubs/sp/800/218/final}. Later drafts are not treated as final requirements.
\item[N3] NIST SP 800-61 Rev. 3, \emph{Incident Response Recommendations and Considerations for Cybersecurity Risk Management} (2025), \url{https://csrc.nist.gov/pubs/sp/800/61/r3/final}.
\item[N4] NIST SP 800-63-4, \emph{Digital Identity Guidelines} (2025), \url{https://csrc.nist.gov/pubs/sp/800/63/4/final}.
\item[N5] NIST, FIPS 203/204/205 (2024), \url{https://csrc.nist.gov/Projects/Post-Quantum-Cryptography/publications}.
\item[O1] OWASP, \emph{API Security Top 10} (2023), \url{https://api-security.owasp.org/editions/2023/en/0x11-t10/}.
\item[C1] CISA, \emph{Known Exploited Vulnerabilities Catalog}, \url{https://www.cisa.gov/known-exploited-vulnerabilities-catalog}.
\item[M1] MITRE ATT\&CK Enterprise, \url{https://attack.mitre.org/matrices/enterprise/}.
\item[G1] GitHub, \emph{REST API endpoints for code scanning}, \url{https://docs.github.com/en/rest/code-scanning}.
\end{description}
\begin{warningbox}\textbf{Interpretation rule.} \handbookref{number} identifies a conceptual anchor in the supplied handbook, not a verbatim quotation or an official requirement. \modernext{} identifies a contemporary extension informed by external references or current engineering practice. \designed{} identifies an independently developed example, template, exercise, metric, or operational procedure. Requirements are binding only if separately adopted through applicable law, contract, policy, or the cited standards; this suite does not assert otherwise.\end{warningbox}

\end{document}
